\section{Universal series and reconstruction}\label{sec:universal}

We use the Chern-zero, virtual-dimension-zero part of Pardon's
semi-Fano universal curve-counting algebra \cite[Theorem 1.2]{Pardon}.
With rational coefficients it is
\[
 H=\QQ[x_{g,m}:g\ge0,\ m\ge1],\qquad
 I=(x_{g,m}),\qquad f_g(z)=\sum_{m\ge0}x_{g,m}z^m,
 \quad x_{g,0}=1.
\]
Each element of $H$ is a polynomial; there is no completion in genus.
The coproduct and division operations are algebra homomorphisms with
\begin{equation}\label{eq:operations}
 \Delta x_{g,m}=\sum_{a+b=m}x_{g,a}\otimes x_{g,b},\qquad
 \rho_d x_{g,m}=\begin{cases}x_{g,m/d}&d\mid m,\\0&d\nmid m.
 \end{cases}
\end{equation}
These are the Chern-zero specializations of
\cite[Definitions 3.9, 3.12, Lemmas 4.12--4.13 and Theorem 7.6]{Pardon};
the specialization also uses vanishing in negative virtual dimension.

Let $\Gamma$ be an additive submonoid of a finite-rank torsion-free
lattice, equipped with an integral degree positive off zero and having
finite sublevel sets. The notation $H\llbracket\Gamma\rrbracket$ means
completion in this degree. A series $A=\sum A_\gamma Q^\gamma$, with
$A_0=1$ and $A_\gamma\in I$ for $\gamma\ne0$, is called compatible if
\begin{equation}\label{eq:compatible}
 \Delta A_\gamma=\sum_{\alpha+\eta=\gamma}A_\alpha\otimes A_\eta,
 \qquad \rho_d A_\gamma=A_{\gamma/d}.
\end{equation}
The last coefficient is zero when $\gamma/d$ does not belong to
$\Gamma$. This condition expresses compatibility with addition and
integer multiplication of cycles.

For a smooth projective Calabi--Yau threefold $V$ with torsion-free
$H_2(V;\ZZ)$, let $\Gamma_V$ be its monoid of effective integral curve
classes. An ample divisor gives the required degree. To see finiteness,
represent a basis of the dual real homology space by smooth two-forms.
Their bounded norms bound every coordinate of an effective cycle by a
constant times its K\"ahler area. Bounded area thus bounds its integral
class. The cycle locus $\mathscr C_\gamma(V)$ is compact and open and
closed in the cycle space. Its characteristic function defines
\[
 U_{V,\gamma}=(V,1_{\mathscr C_\gamma(V)})\in H,
 \qquad \Z_V=\sum_{\gamma\in\Gamma_V}U_{V,\gamma}Q^\gamma.
\]
Pulling back the characteristic function under cycle addition gives
$\sum_{\alpha+\eta=\gamma}1_\alpha\boxtimes1_\eta$, and pullback under
multiplication by $d$ gives $1_{\gamma/d}$. Thus $\Z_V$ is compatible.
The constant term represents the empty cycle.

Write $\lambda_h=\GW_{h,0}$ for connected, unmarked GW evaluation.
Pardon's triangularity result \cite[Lemma 7.7]{Pardon} gives
\begin{equation}\label{eq:detectors}
 \lambda_h(I^2)=0,\qquad \lambda_h(x_{g,1})=0\ (g>h),
 \qquad \lambda_h(x_{h,1})=1.
\end{equation}
For $\gamma\ne0$, $\lambda_h(U_{V,\gamma})=N^V_{h,\gamma}$.
Contracted components attached to a nonconstant connected stable map
are included; wholly contracted connected components are excluded from
the corresponding disconnected series.

\begin{lemma}\label{lem:formal}
Compatible series are determined by their connected evaluations in all
classes and genera. Each has a unique factorization
\begin{equation}\label{eq:factor}
 A=\prod_{\gamma\ne0,\,g\ge0}f_g(Q^\gamma)^{e_{\gamma,g}},
 \qquad e_{\gamma,g}\in\QQ,
\end{equation}
with finitely many nonzero exponents at each class.
\end{lemma}
\begin{proof}
Put $p_{g,m}=[z^m]\log f_g(z)$. These are polynomial coordinates on
$H$, and are primitive. A primitive polynomial in these coordinates is
additive, hence linear over $\QQ$. Moreover
$\rho_dp_{g,m}=p_{g,m/d}$ if $d\mid m$, and zero otherwise. By applying
$\rho_M$ at the largest multiplicity in a finite linear combination,
\begin{equation}\label{eq:primitive}
 \Prim(H)\cap\bigcap_{d>1}\ker\rho_d
   =\bigoplus_{g\ge0}\QQ x_{g,1}.
\end{equation}
At a least-degree class where two compatible series differ, the
difference is primitive and killed by all $\rho_d$, $d>1$. Its least
nonzero genus coefficient is detected by \eqref{eq:detectors}.
This proves reconstruction.

For factorization, each coefficient of $L=\log A$ is primitive.
Write $L_\gamma=\sum a_{\gamma;g,m}p_{g,m}$. Comparing the coefficient
of $p_{g,1}$ in $\rho_mL_\gamma=L_{\gamma/m}$ gives
$a_{\gamma;g,m}=a_{\gamma/m;g,1}$. Thus
\[
 L_\gamma=\sum_{m:\,\gamma/m\in\Gamma}\sum_g
 e_{\gamma/m,g}p_{g,m},\qquad e_{\gamma,g}=a_{\gamma;g,1}.
\]
This proves both existence and uniqueness. Polynomiality and bounded
degree make every displayed coefficient a finite sum.
\end{proof}

We also use the following support consequence. If $F\subset
\Gamma\setminus\{0\}$ is closed under positive integer multiplication,
and $\lambda_h((\log A)_\gamma)=0$ for every $h$ and $\gamma\notin F$,
then $\Supp\log A\subset F$. Indeed, every proper divisor of a class
outside $F$ is outside $F$. Induction on degree, followed by
\eqref{eq:primitive} and \eqref{eq:detectors}, proves the assertion.
Since $\lambda_h$ kills $I^2$, it is equivalent to test $A_\gamma$.

\begin{lemma}\label{lem:pt}
For every $g\ge0$,
\[
 \PT(x_{g,1})=z^{1-g}(1+O(z)).
\]
If $R\in\Prim(H)\cap\bigcap_{d>1}\ker\rho_d$ and
$\PT(R)\in z\QQ[[z]]$, then $R=([z]\PT(R))x_{0,1}$.
\end{lemma}
\begin{proof}
Choose a genus-$g$ curve $C$ and a degree-$g-1$ line bundle $L$ with
$H^0(L)=H^1(L)=0$; for $g=0$ use $\mathcal O(-1)$.
The bundle $L\oplus K_C L^{-1}$ has no cohomology and determinant
$K_C$. Every compact degree-one cycle in its total space is the graph
of a section, hence is $C$. Pardon's definition of $x_{g,1}$ can
therefore be taken with no cut-down sections. Purity forces a curve
scheme with this cycle to be $C$ itself. The minimal stable pair is
$\mathcal O\to\mathcal O_C$, at Euler characteristic $1-g$; its
moduli space is a smooth point, with virtual count $+1$.
This proves the leading term. By \eqref{eq:primitive}, $R$ is a finite
sum of $x_{g,1}$. Its largest genus would contribute a unique lowest
power of $z$. All positive genera therefore vanish under the stated
lower bound.
\end{proof}
